% Conclusion
% Source pages: 233–241

\chapter*{Conclusion}
\addcontentsline{toc}{chapter}{Conclusion}
\markboth{Conclusion}{Conclusion}

\srcpage{233}

Before summing up the foregoing examination of a number of questions connected with spatial
constructions in the visual arts, a few words should be said about the problems that were
deliberately set aside in the main body of the book.

Works conveying a three-dimensional world by the means of two-dimensional art by no means
always existed ``in themselves''; they were frequently parts of more complex ensembles.  In
the art of ancient Egypt one can observe many cases in which a unified composition is made
up of sculpture in the round standing flush against a wall, with reliefs to the left and
right of it.  Ancient Egyptian reliefs and wall paintings existed only as elements of
particular architectural structures---temples, palaces, tombs, and so on.  Exceptions
to this---for example, illustrations to the Book of the Dead---did not define the
character of ancient Egyptian art.  The same is true of medieval Europe, where (setting
aside manuscript miniatures) painting, frescoes, and mosaics formed an essential element
of temple architecture.

This synthetic character must have made itself felt both in the painting and relief, and,
on the other hand, influenced the reception of these works, which were contemplated not as
something self-sufficient but as elements of complex, objectively three-dimensional
structures that also had clearly defined functional characteristics (temple, palace, etc.).
Such an approach, though entirely natural in art-historical analysis, cannot be fruitfully
pursued here due to the complete absence of materials that would allow even individual
aspects of such a problem to be subjected to mathematical analysis---for example, the
peculiarities of the visual perception of images on flat and on curved surfaces, the
influence of the scale and proportions of a building on the perception of images of
various sizes and formats, and so on.

Another question not touched upon above, yet whose importance was indicated in the
introductory chapter, is the question of the history of the development of methods of
spatial construction.  This aspect, requiring deep analysis of the history of the visual
arts, likewise could not become a subject of examination here, especially since the book
studied far from all, but only the most characteristic, types of spatial construction
that emerged in the course of the development of the visual arts.  It is precisely this
circumstance that conditioned the subtitle ``Sketch of the Principal Methods.''

The chapters above attempted to show that the methods of spatial construction
characteristic of different epochs and regions are each reasonable in themselves,
grounded in the objective laws of geometry and the psychology of perception, and cannot
be compared with one another from the standpoint of ``better or worse,'' ``more or less
correct.''  Even Cézanne's perspective system cannot always be considered definitely
better than the classical system of linear perspective, since the latter more faithfully
follows the retinal image.  The approach adopted in the book, which essentially excluded
the analysis of the development of geometric pictorial means, gave the entire exposition
a somewhat static character.

Without claiming any comprehensive treatment of the complex problem of the evolution of
methods of spatial construction---even in its purely geometric aspect---let us offer here
a few reflections on the circle of questions touched upon.

If one arranges the geometric pictorial means that arose in the course of historical
development in a rough chronological order, one can assert that the present book examined
four principles of depicting three-dimensional space on a plane that arose successively:
drafting methods, axonometric methods (and their transformations), the system of classical
linear perspective (discussed in the book in the process of comparison with other methods
and systems), and finally the rigid perceptual perspective system (as exemplified by
Cézanne's landscape painting).  This sequence is broadly comparable to the historical
course of development.  Moreover, it essentially encompasses the whole body of visual art
from early antiquity to modern times.  Indeed, drafting methods are characteristic not
only of Egypt but of many other cultures of antiquity and even the Middle Ages.  Following
the drafting methods there arise methods of conveying perceptual space: first on the basis
of axonometry and its natural transformations (classical and late antiquity, the Middle
Ages from Europe to the Far East), and then, in the Renaissance, by means of the system
of central linear perspective.  This sequence is closed---since the art of the most recent
period is not treated in the book---by the system of central curvilinear perspective.

\srcpage{234}

If one views this sequence from the standpoint of geometry, one can observe the gradual
elaboration of the systems employed, that is, register an element of development.  All four
methods rest on the principle of projection, carried out, however, in different ways and
developed in different ways.

Drafting methods are based on projection from an infinitely remote centre on the condition
that the projection lines are orthogonal to the picture plane.  In the axonometric system
the requirement of orthogonality is dropped and consequently the centre of projection may
be located arbitrarily, but must remain infinitely remote.  In the Renaissance system of
central projection the condition of infinite remoteness of the projection centre is also
dropped.  Thus geometrically each subsequent system possesses greater generality and in a
certain sense greater complexity.  The discussion here concerns, of course, the position
of the projection centre, not the artist's eye; the two points need not coincide.

Projections thus obtained undergo transformations in certain cases.  In medieval art,
axonometry is only a foundation which, through free transformations, gives way to reverse
or (much more rarely) direct perspective; such variations may in part be conditioned by
the laws of natural visual perception.  As for the space of the picture as a whole, it is
constructed on the basis of local axonometries and their modifications.

The Renaissance perspective system claims to convey the entire space of the picture as a
whole; yet its inherent shortcomings long compelled artists to introduce artistically
justified modifications into it.  In Cézanne this general tendency toward refinement of
the linear perspective system led to a qualitative leap---one can no longer speak here of
reasonable ``corrections'' to the Renaissance ways of conveying space, but must recognize
a transition to a new perspective system.  The transition to the rigid perceptual
perspective system that he made intuitively, from a mathematical standpoint, attests to
a further elaboration of the theoretical foundations of perspective construction.  Here
the projection method, analogous to the Renaissance one, is supplemented by a
transformation of this projected image by continuous functions of a certain class, with
the result that the final system cannot be reduced to a projection as the three preceding
ones can.

\srcpage{235}

It would hardly be right to assert that artists consciously recognized this increasing
elaboration of the foundations of spatial construction methods as they developed.  The
increasing mathematical complexity occurred ``by itself'' in connection with the change,
and as a rule the increasing complexity, of the tasks set before artists in the course of
the historical development of culture.

Although the elaboration mentioned above did take place, this does not at all mean that
each subsequent one of the four described stages was, from a mathematical standpoint, more
correct than the previous.  Each of these four stages is, from the standpoint of its
internal logic, mathematically and psychologically irreproachable.  It is precisely this
circumstance that leads to the conclusion that any of these four different methods of
spatial construction is capable of ``revival'' and, if the artistic specifics of a work
require it, may be successfully employed today.

Concluding this brief survey of certain aspects of the historical development of methods
of spatial construction, it seems useful to return once more to the transition from
conveying the geometry of objective space (drafting methods) to conveying the geometry
of perceptual space (axonometry, central linear perspective, central curvilinear
perspective).  While the sequential replacement of the latter three methods is easily
interpreted from the standpoint of evolution and gradual development of pictorial
rendering, the transition from drawing to picture is genuinely revolutionary.

It would be wrong to say that, for example, in ancient Egypt the level of spatial thinking
was not yet high enough and the Egyptians proved unable to ``invent'' even axonometry.
The matter is apparently more complex.  For many examples of ancient Egyptian art contain
drawings so perfect and complex (for instance, the cross-section of a three-storey house
mentioned in Chapter~I, conveying the constructional features of the building, its floors,
staircases, and so on, with remarkable exactness) that their execution requires
considerably more powerful spatial imagination than any picture.  Therefore the transition
from drawing to picture was connected not with a turn to more complex spatial thinking but
with a fundamental change in the orientation of artist and viewer.

\srcpage{236}

The drafting approach is characteristic not only of the art of ancient Egypt but also of
the overwhelming majority of the initial stages of other cultures.  Its formation and
development become understandable if one assumes that in this case the artist conveys not
the visual perception of space (as is often written and thought) but the integral
perception of space.  For this representation is formed not only by the subconscious
analysis of a momentary visual image, but also by movement, tactile contact with the
surfaces of objects, and so on.  Moreover, it is precisely these latter channels for
receiving spatial information that prove to be the more accurate, and the brain
subconsciously approximates the aggregate impressions from many visual images toward
these non-visual but more adequately geometry-conveying representations.  Only a drawing
can convey this general, aggregate representation of space.  Axonometry looks completely
absurd if the visual channel of incoming spatial information is mentally excluded.

Suppose the artist has reproduced objective space in a wall painting or on a sheet of
papyrus.  By virtue of the specifics of pictorial art, this information can only be read
visually.  A kind of paradoxical situation arises: by reproducing objective space, with
the emphasis on \emph{knowledge} rather than \emph{vision}, the artist and viewer can
assimilate this information only visually---that is, by precisely the path which,
conditionally speaking, was considered secondary in natural spatial perception.

In the end, there arises an understandable desire to align the channels for receiving and
transmitting information.  Since the artist can only transmit information visually, it is
the ``input'' as well (to use this perhaps not entirely apt cybernetic term)---that is,
what enters the artist's consciousness---that must consist purely of visual information;
all else is to be ignored, filtered out.  As soon as such a fundamentally new orientation
prevails, the drawing gives way to the picture.

In light of the above, one can assert that the transition from drawing to picture is
connected not with a deepening and development of spatial concepts but with the conscious
severing of a number of channels through which spatial information enters human
consciousness---with a radical change of orientation.  The aggregate use of all the
channels of spatial information available to a person is entirely natural, while the
severing of several such channels is sufficiently arbitrary.  It is for this reason that
drawing precedes picture in the history of art.

\begin{center}$*\quad*\quad*$\end{center}

\srcpage{237}

Let us now turn to the conclusions that follow from the analysis of spatial construction
methods carried out in this book.  Rigorous mathematical analysis has revealed that there
never has existed and never can be developed a scientific system of spatial construction---
in particular a scientific perspective system---that adequately conveys the geometric
characteristics of the depicted space on the picture plane without any conventions or
``distortions.''  In this connection one must acknowledge that even the system of linear
perspective, created by artists of the Renaissance and propagated for centuries, is no
exception to this rule.  Moreover, it cannot be considered the system that most fully
conveys deep space (even when all the restrictions on its use recommended by modern
theory are observed), since, for example, the variant of a mathematical perceptual
perspective system worked out in this book merely as an example---while possessing many of
the advantages of linear perspective---carries approximately half the distortions of
natural visual perception of space that the latter does.  But in that case the fetishization
of the system of linear perspective (unfortunately still current) is not only mistaken but
harmful, since it leads to entirely wrong evaluations of other methods of spatial
construction.

If, consequently, an ideal (mathematical) system of spatial construction in painting
cannot exist, then a range of varieties of scientific spatial construction systems must
inevitably coexist, each best suited for the most faithful possible rendering not of all
but of quite specific geometric features of the subject.  These partial systems will then
not compete with one another but complement each other, since they will have different
domains of application.  Therefore, before evaluating any possible spatial construction,
one should formulate the geometric aspect of the task facing the artist.

Most essential in such a formulation is the answer to the question of which kind of
space---objective or perceptual---is to be conveyed on the picture plane.  In the first
case, drafting methods will prove necessary and a painting of the type given by the art
of ancient Egypt will be realized.  In the second, more common case---when the task of
conveying the geometry of perceptual space is posed---it will be necessary to use
different variants of the scientific perspective system depending on what is to be depicted.
Here it becomes significant whether one is dealing with the depiction of very near and
shallow spaces, or of a space in which the foreground, middle ground, and background are
equally important; whether conveying horizontal or vertical structures is more
important, and so forth.  For each of these tasks a particular variant of mathematically
grounded spatial construction proves preferable.  The matter is further complicated if,
alongside the primary aim of conveying the geometry of perceptual space, the partial
transmission of the geometry of objective space is also possible or even

\srcpage{238}

desirable, as is the case in medieval and modern art.

Into this apparently great variety of scientific spatial construction methods one can
introduce a certain order if they are classified according to what is to be depicted and
how the problem of inevitable distortions is resolved.  One can here speak of three
fundamental methods of depicting space on a plane:

\begin{enumerate}
  \item Drafting methods (conveying the geometry of objective space).
  \item System of perceptual perspective (conveying the geometry of perceptual space):
    \begin{itemize}
      \item[a)] Rigid---in it the question of inevitable distortions is resolved in advance
        with respect to the entire picture space and ``rigidly'' carried through;
      \item[b)] Free---in it the question of inevitable distortions is resolved differently,
        with respect to each individual object depicted in the picture.
    \end{itemize}
\end{enumerate}

If the geometric order of an image using drafting methods is fairly constant, then the
differences in the geometry of perspective constructions are connected with which
specific elements of the image---in resolving the artistic task at hand---it will be
considered admissible to render geometrically distorted, and which not (for an image
entirely free of such distortions is possible only in the uninteresting and trivial case
of conveying the geometry of strongly remote regions of space).

Depending on the choice made, one or another variant of the perceptual perspective
system is realized.  Some examples of such systems were examined in the present book.
It should be emphasized, however, that this apparent multiplicity in fact represents a
strict unity, since all these variants of perspective systems proceed from one and the
same premise (the transformation of the retinal image in the process of visual perception
by constancy mechanisms common to all of them) and are described by one and the same
mathematical relations, in which---speaking figuratively---the emphases are merely placed
differently.

Such unity permits one to speak of there being only one system of scientific perceptual
perspective, with all the above-mentioned variants contained in it as subsystems.  These
subsystems---from medieval reverse perspective to Cézanne---however different they may
appear on the surface, reveal a deep genetic kinship and therefore, in particular, admit
a generalized descriptive characterization: in the system of scientific perceptual
perspective, distant regions of space tend toward ordinary linear perspective, near
regions toward axonometry and weak reverse perspective, and on the middle plane the images
of objectively straight lines have maximum curvature.\footnote{This of course does not
  mean that, for example, modern art differs from medieval art merely in the differing
  remoteness and depth of the ``mastered'' space.}

\srcpage{239}

The subsystems discussed above for the unified system of perceptual perspective were
divided into two groups (depending on the rigid or free choice of distorted elements).
In such a classification, the familiar system of linear perspective turns out to be one
special case of the rigid perceptual system.  Although linear perspective can be regarded
as one of the subsystems of the unified system introduced here, it occupies a special
position among them, since it approximates the retinal image produced by contemplating
the picture to the retinal image produced by nature.  The system of linear perspective
has a dual character: it can be regarded as a special case of perceptual perspective, and
also as a system independent of the latter.  This is due to the fact that, being an
idealized reproduction of the retinal image, it can be interpreted as an image connected
with the first stage of the visual perception process---that is, without taking into
account the activity of the brain---in which case calling it ``perceptual'' is
inappropriate.  Yet simultaneously it also corresponds to the visible image with due
allowance for the transforming activity of the visual perception system in that special
case where this transformation reduces to a ``stretching'' that preserves similarity---
that is, without changing the geometry of the retinal image (which, as has been
repeatedly emphasized, is characteristic of the perception of distant regions of space)---
and in that case it is unquestionably ``perceptual.''  This dual character of the system
of linear perspective led to the situation in the present book where, depending on the
question under discussion, it was sometimes contrasted with the perceptual perspective
system and sometimes included in it as a special case.

Artists are, of course, by no means obliged to follow strictly the laws of visual
perception.  Working within the linear perspective system, a painter has every right to
render comparatively near regions by its laws, just as an icon painter who renders the
near foreground in slight reverse perspective may show the architectural staffage of the
second plane---contrary to the laws of visual perception---also in reverse perspective.
In both these cases the artists will be working within the perspective system characteristic
of them (or of their era), which by no means necessarily corresponds exactly to some
scientific perspective system.

Art historians carrying out the analysis of artworks must always take into account the
multiplicity of scientific perspective systems and, when analyzing an artist's work, must
clearly understand which variant of the perspective system the artist mainly adheres to
and what the deviations from it amount to.  Otherwise such analysis will suffer from
insufficient objectivity.

V.\,N.~Lazarev, in the book \emph{Old Italian Masters}, characterizing the painting of
Duccio, says that he ``commits numerous errors'' in perspective constructions.\footnote{V.\,N.~Lazarev.
  \emph{Old Italian Masters}.  Moscow, 1972, p.\,16.}  This is true if one adds
the words ``with respect to the system of linear perspective,'' which the author
unquestionably has in mind.  If one assumes that Duccio worked in the system of free
perceptual perspective, then his constructions contain almost no errors---that is, his
constructions correspond to the perspective features described in Chapters~IV, V,
and~VI of the present book.  Thus the assessment of an artist's work (whether he erred
or not) depends on how adequate to the artist's work is the scientific perspective the
art historian has in mind.

\srcpage{240}

Since for the near foreground the system of linear perspective gives excessively large
distortions, it is not surprising that even Mantegna---fully conversant with it and a
vigorous champion of it---nevertheless preferred for his
\emph{Dead Christ}\footnote{Ibid., ill.\,147.}
the system of perceptual perspective for near space: axonometry.  This can easily be
verified by comparing the width of Christ's feet in the picture with the distance between
his eyes.  The absence in the image of any diminution of Christ's body with increasing
depth indicates that we are here dealing with perceptual perspective.  By depicting the
wooden bier on which the body of Christ lies as narrowing in depth, Mantegna made an
error in the perspective construction relative to the system of perceptual perspective
that governs his picture.  Of course, the word ``error'' carries no evaluative
character here.  If one were to assert, starting from the usual fetishization of the
system of linear perspective, that Mantegna painted the bier ``correctly'' and the body
of Christ ``in error,'' this would give a distorted conception of the artist, who is
allegedly erring in the principal thing.

F.~Novotny, in his monograph widely cited in Chapter~X, finds a multitude of errors in
Cézanne, holding that the only scientific perspective is linear perspective.  If one
recognizes that the rigid perceptual system is no less scientific, the judgment about the
features of spatial construction in Cézanne's paintings turns out to be directly
opposite---this artist makes almost no deviations from the scientific perspective system
that he intuitively followed.  Furthermore, if one analyses Cézanne's painting by
relying on the theory of linear perspective, then, as was shown in Chapter~X, the art
historian finds himself compelled to attribute very high viewpoints to Cézanne even in
cases where the artist did not employ them.

These several examples clearly show the necessity of understanding which variety of
scientific perspective the artist favors, since a lack of clarity on such a question
can lead to entirely distorted conceptions of his work.  Appendix~8 gives a table of
function values and the formulas determining the principal types of rigid perceptual
perspective systems; as for other variants of such systems, they are determined by
analogous expressions and therefore cannot yield anything qualitatively new.  The
general character of such systems is well enough illustrated in Chapter~X.

\srcpage{241}

The fetishization of the system of linear perspective in the study of painting that
conveys the illusion of spatial depth can, as was shown above, lead to erroneous
evaluations of the work of artists.  Even more inappropriate is the examination of
painting that conveys objective space (for example, in the art of ancient Egypt) from
the same standpoint of classical perspective, since no perspective system has the
slightest relation to drafting methods of representing space on the picture plane.
Attempts to ``justify'' ancient Egyptian painting by appealing to the naivety of the
artist or his adherence to traditional methods of depiction should be considered
untenable, since both these approaches implicitly assume that the Egyptians had not
``grown up to'' knowing perspective.  Yet in ancient Egypt the tradition consisted not
so much in a set of pictorial methods as in the requirement to convey the geometry of
objective space, and this geometry was conveyed in an optimal manner.  Consequently
the method of depiction itself---the type of its geometry---was conditioned by objective
reasons.  To traditional artistic features one can attribute only the choice of specific
drafting conventions, that is, features of a secondary character.  Whenever the
requirement to convey objective space arose anywhere, the Egyptian method of depiction
spontaneously reappeared, demonstrating that tradition plays a subordinate role here.
This is attested to, in particular, by the example of the Indian (pre-Mughal) miniature
adduced in Chapter~IX.

Medieval perspective painting includes, as coequal elements, depictions of objects
rendered by drafting methods.  This complicates the geometric structure of medieval
pictorial art, and this feature must always be borne in mind in artistic analysis.

The recourse to geometrically contradictory images, likewise characteristic of the Middle
Ages, is grounded in the objective properties of human psychology.  As was shown in
Chapters~VIII and~IX, this device---deployed with great tact---frequently served to
intensify the expressiveness of a pictorial work.  Yet today in the artistic analysis of
compositions where such a device is employed, this geometric feature is as a rule not
examined.  A kind of ``conspiracy of silence'' has formed around geometrically
contradictory images.  Yet, for example, the geometry of Rublev's
\emph{Trinity} and \emph{Presentation} cannot be understood without recourse to the
analysis of geometric self-contradiction.  Therefore an art-historical analysis of this
intelligent feature of medieval painting would likely yield much of interest.

As the examination of typical methods of spatial construction carried out in this book
has shown, all of them have a rational foundation; and it is precisely this that permits
the application to their analysis of the exact methods of natural-scientific inquiry---
methods that cannot, of course, claim to disclose the artistic image, but may prove
entirely appropriate in attempts to arrive at a deeper understanding of the geometry of
pictorial means.

